\documentclass[12pt,letterpaper]{article}

% --- PAQUETES NECESARIOS ---
\usepackage[utf8]{inputenc}
\usepackage[spanish,es-tabla]{babel} % Idioma español
\usepackage[margin=2.5cm]{geometry}  % Márgenes
\usepackage{graphicx}                % Para insertar imágenes
\usepackage{amsmath,amssymb}         % Símbolos matemáticos
\usepackage{setspace}                % Interlineado
\usepackage{hyperref}                % Enlaces interactivos (sin guion)
\usepackage{cite}                    % Manejo de citas bibliográficas

\onehalfspacing % Interlineado de 1.5

\begin{document}

% =========================================================
% PORTADA INSTITUCIONAL
% =========================================================
\begin{titlepage}
    \centering
    
    % --- ENCABEZADO CON LOGOS ---
    % Reemplaza "logo_tecnm.png" y "logo_cenidet.png" con los nombres reales de tus imágenes en Overleaf
    \begin{minipage}{0.45\textwidth}
        \begin{flushleft}
            \includegraphics[height=1.8cm]{example-image-a} % Imagen simulada para TecNM
        \end{flushleft}
    \end{minipage}
    \hfill
    \begin{minipage}{0.45\textwidth}
        \begin{flushright}
            \includegraphics[height=1.8cm]{example-image-b} % Imagen simulada para CENIDET
        \end{flushright}
    \end{minipage}
    
    \vspace{1cm}
    
    % --- INSTITUCIÓN ---
    {\bfseries\large TECNOLÓGICO NACIONAL DE MÉXICO} \\[0.3cm]
    {\bfseries\large INSTITUTO TECNOLÓGICO DE CUAUTLA} \\[0.8cm]
    
    % --- PROGRAMA Y SEMESTRE ---
    {\bfseries\large INGENIERÍA EN SISTEMAS COMPUTACIONALES} \\[0.2cm]
    {\bfseries\large ESPECIALIDAD INTELIGENCIA DE NEGOCIOS Y CIENCIA DE DATOS} \\[0.2cm]
    {\bfseries\large SÉPTIMO SEMESTRE} \\[0.8cm]
    
    % --- PERIODO Y MATERIA ---
    {\bfseries AGOSTO - DICIEMBRE 2026} \\[0.2cm]
    {\bfseries INFORME TÉCNICO - UNIDAD I} \\[0.8cm]
    
    % --- TÍTULO DE LA ACTIVIDAD / PROYECTO ---
    {\bfseries\Large PROYECTO INTEGRADOR 0} \\[0.3cm]
    {\bfseries\large DISEÑO Y SIMULACIÓN DE SISTEMAS CONTROLADOS} \\[1.2cm]
    
    % --- PRESENTADO POR ---
    {\small PRESENTADO POR:} \\[0.1cm]
    {\bfseries ARÍSTIDES CABALLERO ALFARO} \\[0.6cm]
    
    % --- DIRECTOR ---
    {\small DOCENTE:} \\[0.1cm]
    {\bfseries ARÍSTIDES CABALLERO ALFARO} \\[0.6cm]
    
        
    % --- LUGAR Y FECHA ---
    \vfill
    {\bfseries CUAUTLA, MORELOS, MÉXICO} \\
    {\bfseries SEPTIEMBRE, 2026}

\end{titlepage}

% =========================================================
% ÍNDICES
% =========================================================
\newpage
\pagenumbering{roman} % Numeración romana para los índices

\tableofcontents % Índice General
\newpage

\listoffigures % Índice de Figuras
\newpage

% =========================================================
% CUERPO DEL DOCUMENTO
% =========================================================
\pagenumbering{arabic} % Numeración arábiga a partir de la Introducción

\section{Introducción}
Esta es una sección de prueba para la introducción del reporte. En esta sección se presenta una visión general del tema a desarrollar, contextualizando el problema y la justificación del trabajo planteado.

Como parte del desarrollo, se pueden agregar referencias bibliográficas para sustentar el marco teórico y metodológico, por ejemplo: este trabajo se basa en las metodologías presentadas por \cite{mendoza2021} y los conceptos fundamentales expuestos en \cite{smith2020}.

\section{Objetivos de la actividad}
\subsection{Objetivo General}
Diseñar una metodología para el desarrollo de juegos formativos orientados al aprendizaje significativo mediante un enfoque constructivista.

\subsection{Objetivos Específicos}
\begin{itemize}
    \item Analizar los principios del enfoque didáctico constructivista aplicados en entornos virtuales.
    \item Implementar modelos de sistemas retroalimentados para la evaluación continua.
    \item Validar el impacto de la herramienta mediante un caso de estudio práctico.
\end{itemize}

\section{Marco Teórico}
En esta sección se describen los fundamentos teóricos necesarios para comprender el desarrollo de la investigación o reporte.

\subsection{Enfoque Constructivista}
El constructivismo sostiene que el aprendizaje es un proceso activo donde el estudiante edifica nuevas ideas o conceptos basándose en su conocimiento actual y pasado.

\subsection{Sistemas Retroalimentados}
Un sistema retroalimentado permite ajustar la salida en función de las respuestas o resultados previos, garantizando un flujo continuo de mejora.

\section{Desarrollo}
En esta sección el estudiante debe describir los pasos, procedimientos, algoritmos o experimentos realizados durante la actividad.

A continuación, se presenta un ejemplo de cómo incluir una figura en el documento para que aparezca automáticamente en el \textbf{Índice de Figuras}:

\begin{figure}[htbp]
    \centering
    \includegraphics[width=0.5\textwidth]{example-image} % Imagen simulada
    \caption{Diagrama de bloques del sistema retroalimentado.}
    \label{fig:diagrama_sistema}
\end{figure}

Como se observa en la Figura \ref{fig:diagrama_sistema}, el flujo de información se mantiene continuo a lo largo del proceso.

\section{Conclusión}
En las conclusiones se condensan los hallazgos principales alcanzados durante el desarrollo del trabajo, evaluando el cumplimiento de los objetivos planteados y proyectando posibles trabajos futuros.

% =========================================================
% BIBLIOGRAFÍA
% =========================================================
\newpage
\bibliographystyle{ieeetr} % Estilo IEEE (puedes cambiarlo por apa, plain, etc.)
\bibliography{referencias} % Archivo .bib

\end{document}